\documentclass[12pt,a4paper]{article}

% --- PAQUETES DE IDIOMA Y FORMATO ---
\usepackage[spanish,es-tabla]{babel}
\usepackage[utf8]{utf8}
\usepackage[margin=2.5cm]{geometry}
\usepackage{helvet}
\renewcommand{\familydefault}{\sfdefault}

% --- COLORES Y ESTILOS ---
\usepackage{xcolor}
\definecolor{primaryColor}{RGB}{15, 23, 42}       % Slate 900
\definecolor{secondaryColor}{RGB}{30, 58, 138}    % Blue 900
\definecolor{accentColor}{RGB}{16, 185, 129}      % Emerald 500
\definecolor{codeBg}{RGB}{248, 250, 252}          % Slate 50
\definecolor{codeBorder}{RGB}{226, 232, 240}      % Slate 200
\definecolor{codeKeyword}{RGB}{147, 51, 234}     % Purple 600
\definecolor{codeString}{RGB}{5, 150, 105}        % Emerald 600
\definecolor{codeComment}{RGB}{100, 116, 139}     % Slate 500

% --- PAQUETES GRÁFICOS Y DE CÓDIGO ---
\usepackage{listings}
\usepackage{tcolorbox}
\usepackage{hyperref}
\usepackage{amsmath}
\usepackage{amssymb}
\usepackage{fancyhdr}
\usepackage{tikz}
\usetikzlibrary{shapes.geometric, arrows, positioning}

% --- CONFIGURACIÓN DE LISTINGS (CÓDIGO FUENTE) ---
\lstdefinelanguage{TypeScript}{
  keywords={typeof, new, true, false, catch, function, return, null, catch, switch, var, if, in, while, do, else, case, break, export, import, from, class, interface, extends, implements, constructor, private, public, protected, readonly, inject, signal, computed, Observable},
  keywordstyle=\color{codeKeyword}\bfseries,
  ndkeywords={class, export, boolean, number, string, void, any, Injectable, Component},
  ndkeywordstyle=\color{secondaryColor}\bfseries,
  identifierstyle=\color{primaryColor},
  sensitive=true,
  comment=[l]{//},
  morecomment=[s]{/*}{*/},
  commentstyle=\color{codeComment}\itshape,
  stringstyle=\color{codeString},
  morestring=[b]',
  morestring=[b]"
}

\lstset{
  language=TypeScript,
  backgroundcolor=\color{codeBg},
  basicstyle=\ttfamily\small,
  breaklines=true,
  captionpos=b,
  frame=single,
  rulecolor=\color{codeBorder},
  tabsize=2,
  showstringspaces=false
}

% --- ENCABEZADO Y PIE DE PÁGINA ---
\pagestyle{fancy}
\fancyhf{}
\rhead{\textcolor{secondaryColor}{\textbf{Gobernación del Cauca --- Sistema de Liquidadores}}}
\lhead{\textcolor{codeComment}{Documentación de Infraestructura HTTP Core}}
\rfoot{Página \thepage}

\hypersetup{
    colorlinks=true,
    linkcolor=secondaryColor,
    urlcolor=accentColor,
    pdftitle={Documentación Técnica: Infraestructura HTTP y Consumo de APIs en Angular},
    pdfauthor={Equipo de Desarrollo --- Gobernación del Cauca}
}

\begin{document}

% --- PORTADA ---
\begin{titlepage}
    \centering
    \vspace*{1cm}
    
    {\Huge \bfseries \textcolor{primaryColor}{SISTEMA DE LIQUIDADORES}}\\[0.4cm]
    {\Large \textcolor{secondaryColor}{GOBERNACIÓN DEL CAUCA}}\\[1.5cm]
    
    \rule{\linewidth}{0.5mm}\\[0.4cm]
    {\huge \bfseries \textcolor{primaryColor}{Infraestructura HTTP Genérica,\\ Interceptores y Consumo por Capas}}\\[0.2cm]
    \rule{\linewidth}{0.5mm}\\[1.5cm]
    
    \begin{tcolorbox}[colframe=secondaryColor,colback=codeBg,arc=3mm,title=\textbf{Resumen Ejecutivo}]
        Documento técnico de especificación sobre el diseño, arquitectura e implementación del módulo \textbf{Core} en Frontend (Angular 21+). Se detalla la función HTTP genérica con soporte multi-impuesto, la gestión reactiva con Angular Signals, los interceptores de seguridad JWT y el patrón de consumo por capas (Domain $\rightarrow$ Infrastructure $\rightarrow$ Application $\rightarrow$ Presentation).
    \end{tcolorbox}
    
    \vfill
    
    {\large \textbf{Fecha:} Agosto 2026}\\[0.2cm]
    {\large \textbf{Tecnologías:} Angular 21+, RxJS, TypeScript, ASP.NET Core API}\\[0.2cm]
    {\large \textbf{Estado:} Implementado y Verificado}
    
    \vspace*{1cm}
\end{titlepage}

\tableofcontents
\newpage

% --- SECCIÓN 1 ---
\section{Introducción y Arquitectura}

El \textbf{Sistema de Liquidadores de la Gobernación del Cauca} está fundamentado en una arquitectura distribuida y modularizada. A nivel de Backend, el sistema separa la responsabilidad tributaria en múltiples APIs Web independientes:

\begin{itemize}
    \item \textbf{Login / Orquestador API:} Gestiona la autenticación centralizada, emite tokens JWT y devuelve la información del módulo tributario asignado junto a su URL base.
    \item \textbf{API Impuesto de Automotores:} Modulo encargado de vehículos, propietarios, liquidaciones y tarifas de automotores.
    \item \textbf{API Impuesto de Registros:} Módulo dedicado a actos, sujetos y liquidaciones del impuesto de registro.
\end{itemize}

Para reflejar este principio en el Frontend (Angular), la carpeta \texttt{src/app/core/} aloja los componentes transversales que gestionan de forma dinámica la comunicación HTTP, el almacenamiento de tokens y la seguridad global.

---

% --- SECCIÓN 2 ---
\section{Infraestructura del Core HTTP}

\subsection{Modelo Matemático de Enrutamiento de URLs}

Cada impuesto posee su propia dirección IP o dominio. La construcción final de una petición HTTP se representa formalmente como la concatenación de la Base URL del módulo objetivo y la ruta relativa del recurso:

\[
\text{URL}_{\text{final}}(m, e) = \text{BaseUrl}(m) \mathbin{\Vert} \text{CleanEndpoint}(e)
\]

Donde:
\begin{itemize}
    \item $m \in \{\text{AUTOMOTORES}, \text{REGISTROS}, \text{LOGIN}, \dots\}$ es el módulo o impuesto seleccionado.
    \item $e$ representa la ruta del recurso final (ej. \texttt{vehiculos/ABC123}).
    \item $\text{BaseUrl}(m)$ obtiene la URL base configurada o retornada dinámicamente por la Login API.
\end{itemize}

---

\subsection{Gestión de Estado Reactivo con Signals (\texttt{AuthStateService})}

El servicio \texttt{AuthStateService} administra el estado del usuario autenticado, el módulo activo y la relación de URLs base mediante **Angular Signals**.

\begin{lstlisting}[caption={Fragmento del AuthStateService (Signals reactivos)}]
@Injectable({ providedIn: 'root' })
export class AuthStateService {
  readonly currentUser = signal<User | null>(null);
  readonly currentModulo = signal<TaxModuleType | null>(null);

  readonly moduleApiUrls = signal<Record<string, string>>({
    LOGIN: 'http://localhost:5000/api',
    AUTOMOTORES: 'http://localhost:5001/api',
    REGISTROS: 'http://localhost:5002/api',
  });

  getApiUrl(modulo?: TaxModuleType): string {
    const target = (modulo || this.currentModulo() || 'LOGIN').toUpperCase();
    return this.moduleApiUrls()[target] || '';
  }
}
\end{lstlisting}

---

\subsection{Servicio HTTP Genérico Multi-Impuesto (\texttt{BaseApiService})}

\texttt{BaseApiService} es el núcleo de la comunicación HTTP. Encapsula las operaciones \texttt{GET}, \texttt{POST}, \texttt{PUT}, \texttt{PATCH} y \texttt{DELETE}, abstrayendo la complejidad del armado de URLs y aprovechando los interceptores globales de Angular.

\begin{lstlisting}[caption={Servicio HTTP Genérico BaseApiService}]
@Injectable({ providedIn: 'root' })
export class BaseApiService {
  private http = inject(HttpClient);
  private authState = inject(AuthStateService);

  public buildUrl(endpoint: string, targetModule?: TaxModuleType): string {
    const baseUrl = this.authState.getApiUrl(targetModule);
    const cleanBase = baseUrl.endsWith('/') ? baseUrl.slice(0, -1) : baseUrl;
    const cleanEndpoint = endpoint.startsWith('/') ? endpoint : `/${endpoint}`;
    return `${cleanBase}${cleanEndpoint}`;
  }

  get<T>(endpoint: string, options?: HttpOptions, targetModule?: TaxModuleType): Observable<T> {
    const url = this.buildUrl(endpoint, targetModule);
    return this.http.get<T>(url, this.formatOptions(options));
  }

  post<T>(endpoint: string, body: any, options?: HttpOptions, targetModule?: TaxModuleType): Observable<T> {
    const url = this.buildUrl(endpoint, targetModule);
    return this.http.post<T>(url, body, this.formatOptions(options));
  }
}
\end{lstlisting}

---

\subsection{Security y Interceptores HTTP}

Se implementaron interceptores HTTP funcionales (\texttt{HttpInterceptorFn}) registrados en \texttt{app.config.ts}:

\begin{enumerate}
    \item \textbf{AuthInterceptor:} Extrae el \texttt{accessToken} guardado en \texttt{TokenStorageService} e inyecta la cabecera:
    \[
    \text{Authorization: Bearer } \langle \text{accessToken} \rangle
    \]
    \item \textbf{ErrorInterceptor:} Captura excepciones de red como HTTP 401 (\textit{Unauthorized}) y limpia automáticamente la sesión activa redirigiendo al usuario hacia \texttt{/login}.
\end{enumerate}

---

% --- SECCIÓN 3 ---
\section{Patrón de Consumo por Capas (Features)}

Para preservar la arquitectura limpia en el frontend, el consumo de las APIs desde un módulo (ej. \texttt{automotores}) sigue una separación rigurosa de responsabilidades:

\begin{center}
\begin{tikzpicture}[node distance=1.5cm, auto]
    \node [draw, rectangle, fill=codeBg, rounded corners, width=4cm, align=center] (pres) {\textbf{Presentation}\\ (automotores-index.ts)};
    \node [draw, rectangle, fill=codeBg, rounded corners, width=4cm, align=center, right=1.5cm of pres] (app) {\textbf{Application / Facade}\\ (vehiculos.facade.ts)};
    \node [draw, rectangle, fill=codeBg, rounded corners, width=4cm, align=center, right=1.5cm of app] (infra) {\textbf{Infrastructure API}\\ (vehiculos-api.service.ts)};
    \node [draw, rectangle, fill=codeBg, rounded corners, width=4cm, align=center, right=1.5cm of infra] (core) {\textbf{Core HTTP}\\ (base-api.service.ts)};

    \draw[->, thick, secondaryColor] (pres) -- node[above] {\small Dispara evento} (app);
    \draw[->, thick, secondaryColor] (app) -- node[above] {\small Solicita caso uso} (infra);
    \draw[->, thick, secondaryColor] (infra) -- node[above] {\small Llama genérico} (core);
\end{tikzpicture}
\end{center}

\subsection{Capa 1: Infraestructura (\texttt{VehiculosApiService})}

Clase encargada de conocer los endpoints de su propio módulo llamando al servicio genérico de \texttt{core}:

\begin{lstlisting}[caption={Infraestructura: VehiculosApiService}]
@Injectable({ providedIn: 'root' })
export class VehiculosApiService {
  private api = inject(BaseApiService);

  obtenerPorPlaca(placa: string): Observable<Vehiculo> {
    return this.api.get<Vehiculo>(`vehiculos/${placa}`, {}, 'AUTOMOTORES');
  }
}
\end{lstlisting}

\subsection{Capa 2: Aplicación / Facade (\texttt{VehiculosFacade})}

Gestor de estado reactivo específico del Feature mediante **Signals**:

\begin{lstlisting}[caption={Aplicación: VehiculosFacade}]
@Injectable({ providedIn: 'root' })
export class VehiculosFacade {
  private api = inject(VehiculosApiService);

  readonly vehiculo = signal<Vehiculo | null>(null);
  readonly loading = signal<boolean>(false);
  readonly error = signal<string | null>(null);

  buscarVehiculo(placa: string): void {
    this.loading.set(true);
    this.error.set(null);

    this.api.obtenerPorPlaca(placa).subscribe({
      next: (data) => {
        this.vehiculo.set(data);
        this.loading.set(false);
      },
      error: () => {
        this.error.set('No se encontró el vehículo');
        this.loading.set(false);
      }
    });
  }
}
\end{lstlisting}

\subsection{Capa 3: Presentación (\texttt{AutomotoresIndex})}

Componente de interfaz gráfica que se conecta a la \texttt{Facade} sin invocar peticiones HTTP directamente:

\begin{lstlisting}[caption={Presentación: AutomotoresIndex}]
@Component({
  selector: 'app-automotores-index',
  standalone: true,
  templateUrl: './automotores-index.html',
})
export class AutomotoresIndex {
  readonly facade = inject(VehiculosFacade);
  readonly placaInput = signal<string>('');

  onBuscar(): void {
    if (this.placaInput().trim()) {
      this.facade.buscarVehiculo(this.placaInput());
    }
  }
}
\end{lstlisting}

---

% --- SECCIÓN 4 ---
\section{Conclusiones}

\begin{enumerate}
    \item \textbf{Desacoplamiento total:} Las vistas y componentes no conocen URLs base ni encabezados de autorización.
    \item \textbf{Seguridad centralizada:} El token JWT se inyecta y procesa en un único punto (\texttt{AuthInterceptor}).
    \item \textbf{Escalabilidad:} La integración de nuevos impuestos (ej. Licores, Degüello) no requiere alterar la infraestructura existente en \texttt{core}.
\end{enumerate}

\end{document}
